% !TEX root = master.tex
% !TEX program = lualatex
\lecture{7}{02-04-2026}{Applications of the Residue Theorem}{}

\section{Applications of Residues to Real Integrals}

The residue theorem can be used to compute integrals over simple closed curves, and sometimes even to compute \emph{real} integrals by replacing them with complex integrals. The strategy is always the same:
\begin{itemize}
\item Extend the real integrand to a complex function $f(z)$.
\item Choose a closed contour $\gamma$ that includes the real axis as part of its boundary.
\item Apply the residue theorem to compute $\int_\gamma f(z)\,dz$.
\item Show that the contribution of the non-real part of the contour vanishes in the limit.
\item Recover the desired real integral.
\end{itemize}

\subsection{A motivating example}

Let us compute
\[
\int_{-\infty}^{\infty} \frac{1}{1+x^2}\,dx.
\]
We already know from real analysis that this integral equals $\pi$ (since $\arctan$ is a primitive of $1/(1+x^2)$). Let us recover this result using complex analysis.

Set
\[
f(z)=\frac{1}{1+z^2}.
\]
Given $R>0$, consider:
\begin{itemize}
\item the segment $L_R:[-R,R]\to\mathbb{C}$, $L_R(t)=t$;
\item the upper semicircle $C_R:[0,\pi]\to\mathbb{C}$, $C_R(t)=Re^{it}$.
\end{itemize}
Let $\gamma_R=L_R\cup C_R$ be the closed counterclockwise oriented contour.

The function $f$ has poles of order $1$ at $z=\pm i$. Only $z=i$ lies in the upper half-plane, so it is the only pole inside $\gamma_R$ (provided $R>1$). By the residue theorem,
\[
\int_{\gamma_R} f(z)\,dz = 2\pi i\,\operatorname{Res}_i(f).
\]
The residue is computed via
\[
\operatorname{Res}_i(f)
=\lim_{z\to i}(z-i)f(z)
=\lim_{z\to i}\frac{1}{z+i}
=\frac{1}{2i}.
\]
Hence
\[
\int_{\gamma_R} f(z)\,dz = 2\pi i\cdot \frac{1}{2i}=\pi.
\]
On the other hand, splitting the contour:
\[
\int_{\gamma_R} f(z)\,dz
=\int_{L_R} f(z)\,dz + \int_{C_R} f(z)\,dz
=\int_{-R}^R \frac{1}{1+t^2}\,dt + \int_{C_R} f(z)\,dz.
\]

It remains to show that the contribution of $C_R$ vanishes as $R\to\infty$. We estimate:
\[
\left|\int_{C_R} f(z)\,dz\right|
=\left|\int_0^\pi \frac{1}{1+R^2 e^{2it}}\,iRe^{it}\,dt\right|
\leq \int_0^\pi \frac{R}{R^2-1}\,dt
=\frac{\pi R}{R^2-1}\xrightarrow[R\to\infty]{}0.
\]
Therefore,
\[
\int_{-\infty}^\infty \frac{1}{1+x^2}\,dx = \pi.
\]

\begin{remark}
This method works for any integrand that decays \emph{strictly faster} than $1/R$ as $R\to\infty$, which ensures the integral over the arc vanishes.
\end{remark}

\subsection{Integrals of rational functions times \texorpdfstring{$e^{iax}$}{exp(iax)}}

A more interesting class of integrals concerns functions of the form $R(x)\cos(ax)$ or $R(x)\sin(ax)$, where $R$ is a rational function. These appear naturally in Fourier analysis, in the study of resonances in damped systems, etc.

\begin{theorem}[Real integrals via residues]
Let $R(x)=\dfrac{P(x)}{Q(x)}$ where $P,Q$ are real polynomials such that:
\begin{enumerate}
\item $\deg Q \geq \deg P + 2$;
\item $Q(x)\neq 0$ for all $x\in\mathbb{R}$.
\end{enumerate}
Let $a\geq 0$ and define $f(z)=R(z)e^{iaz}$. Then
\[
\int_{-\infty}^{\infty} R(x)e^{iax}\,dx
=2\pi i \sum_{k} \operatorname{Res}_{z_k}(f),
\]
where the sum is taken over the poles $z_k$ of $f$ in the \emph{upper half-plane} $\{\Im z>0\}$.
\end{theorem}

\paragraph{Observations}
\begin{itemize}
\item The condition $\deg Q\geq \deg P+2$ forces $\deg Q$ to be even (otherwise $Q$ would have a real root).
\item Since $Q$ has real coefficients, its roots come in pairs of complex conjugates. None of them are real by assumption, so half of them lie in the upper half-plane and the other half in the lower half-plane.
\item Using Euler's formula $e^{iax}=\cos(ax)+i\sin(ax)$, the real and imaginary parts of the integral give
\[
\int_{-\infty}^\infty R(x)\cos(ax)\,dx = \Re\!\left(\int_{-\infty}^\infty R(x)e^{iax}\,dx\right),
\]
\[
\int_{-\infty}^\infty R(x)\sin(ax)\,dx = \Im\!\left(\int_{-\infty}^\infty R(x)e^{iax}\,dx\right).
\]
\end{itemize}

\subsubsection*{Why the contour integral over the arc vanishes}

The key estimate is the following: for $z=x+iy$ with $y\geq 0$ and $a\geq 0$,
\[
|e^{iaz}|=|e^{ia(x+iy)}|=|e^{iax}|\cdot|e^{-ay}|=e^{-ay}\leq 1.
\]

\begin{remark}
Note that $|e^{iaz}|$ is \textbf{not} bounded on all of $\mathbb{C}$: in the lower half-plane ($y<0$), the term $e^{-ay}$ blows up. This is precisely why we must close the contour in the \emph{upper} half-plane when $a\geq 0$. For $a<0$, one would close the contour in the lower half-plane instead.
\end{remark}

For $z=Re^{i\theta}$ on $C_R$ with $\theta\in[0,\pi]$, we have $\Im z = R\sin\theta\geq 0$, hence
\[
|e^{iaRe^{i\theta}}|\leq 1.
\]

For the rational part, write
\[
P(z) = b_m z^m + \cdots, \qquad Q(z)=c_n z^n + \cdots, \qquad n\geq m+2.
\]
Then for large $|z|$,
\[
|R(z)| = \left|\frac{P(z)}{Q(z)}\right| \leq \frac{C}{|z|^{n-m}}\leq \frac{C}{|z|^2}.
\]
Therefore
\[
\left|\int_{C_R} f(z)\,dz\right|
\leq \int_0^\pi |R(Re^{i\theta})|\cdot|e^{iaRe^{i\theta}}|\cdot R\,d\theta
\leq \int_0^\pi \frac{C}{R^{n-m}}\cdot 1 \cdot R\,d\theta
=\frac{\pi C}{R^{n-m-1}}\xrightarrow[R\to\infty]{} 0.
\]

\paragraph{Benefit}
We compute the residues of $f$ by finding its poles and their orders, without ever computing a Laurent series explicitly.

\subsection{A worked example}

We compute
\[
I=\int_{-\infty}^{\infty} \frac{x^2}{16+x^4}\cos x\,dx.
\]
This fits the previous setup with $a=1$, $P(x)=x^2$, $Q(x)=16+x^4$. Note $\deg Q=4\geq \deg P+2=4$, and $16+x^4>0$ for all real $x$. By Euler's formula,
\[
I=\Re\!\left(\int_{-\infty}^{\infty} \frac{x^2}{16+x^4}e^{ix}\,dx\right).
\]
Define
\[
f(z)=\frac{z^2}{16+z^4}e^{iz}.
\]

\paragraph{Finding the poles.}
The poles of $f$ are the roots of $16+z^4=0$:
\[
z^4=-16=16e^{i(\pi+2k\pi)}, \qquad k\in\{0,1,2,3\}.
\]
Therefore
\[
z_k = 2 e^{i\left(\frac{\pi}{4}+\frac{k\pi}{2}\right)}, \qquad k=0,1,2,3.
\]
Explicitly:
\[
z_1 = \sqrt{2}+i\sqrt{2}, \quad
z_2 = -\sqrt{2}+i\sqrt{2}, \quad
z_3 = -\sqrt{2}-i\sqrt{2}, \quad
z_4 = \sqrt{2}-i\sqrt{2}.
\]
Only $z_1$ and $z_2$ lie in the upper half-plane, hence only these contribute.

\paragraph{Computing the residues.}
Each $z_k$ is a simple pole. Using the standard trick: if $g(z)=h(z)/Q(z)$ where $Q$ has a simple zero at $z_k$, then
\[
\operatorname{Res}_{z_k}(g)=\frac{h(z_k)}{Q'(z_k)}.
\]
Here $h(z)=z^2 e^{iz}$ and $Q'(z)=4z^3$, so:
\[
\operatorname{Res}_{z_k}(f)
=\frac{z_k^2 e^{iz_k}}{4 z_k^3}
=\frac{e^{iz_k}}{4 z_k}.
\]

For $z_1=\sqrt{2}+i\sqrt{2}=2e^{i\pi/4}$:
\[
\frac{1}{4z_1}=\frac{1}{8 e^{i\pi/4}}=\frac{e^{-i\pi/4}}{8}=\frac{1-i}{8\sqrt{2}}.
\]
Also $e^{iz_1}=e^{i(\sqrt{2}+i\sqrt{2})}=e^{-\sqrt{2}}e^{i\sqrt{2}}$. Therefore
\[
\operatorname{Res}_{z_1}(f)
=\frac{e^{-\sqrt{2}}e^{i\sqrt{2}}(1-i)}{8\sqrt{2}}.
\]

For $z_2=-\sqrt{2}+i\sqrt{2}=2e^{i3\pi/4}$:
\[
\frac{1}{4z_2}=\frac{e^{-i3\pi/4}}{8}=\frac{-1-i}{8\sqrt{2}}.
\]
Also $e^{iz_2}=e^{-\sqrt{2}}e^{-i\sqrt{2}}$. Therefore
\[
\operatorname{Res}_{z_2}(f)
=\frac{e^{-\sqrt{2}}e^{-i\sqrt{2}}(-1-i)}{8\sqrt{2}}.
\]

\paragraph{Summing.}
\[
\operatorname{Res}_{z_1}(f)+\operatorname{Res}_{z_2}(f)
=\frac{e^{-\sqrt{2}}}{8\sqrt{2}}\left[e^{i\sqrt{2}}(1-i)+e^{-i\sqrt{2}}(-1-i)\right].
\]
Expanding using $e^{\pm i\sqrt{2}}=\cos\sqrt{2}\pm i\sin\sqrt{2}$ and collecting terms:
\[
e^{i\sqrt{2}}(1-i)+e^{-i\sqrt{2}}(-1-i)
=2i\sin\sqrt{2}-2i\cos\sqrt{2}
=2i(\sin\sqrt{2}-\cos\sqrt{2}).
\]
Hence
\[
\operatorname{Res}_{z_1}(f)+\operatorname{Res}_{z_2}(f)
=\frac{e^{-\sqrt{2}}}{8\sqrt{2}}\cdot 2i(\sin\sqrt{2}-\cos\sqrt{2})
=\frac{i\,e^{-\sqrt{2}}}{4\sqrt{2}}(\sin\sqrt{2}-\cos\sqrt{2}).
\]

\paragraph{Final result.}
Multiplying by $2\pi i$:
\[
2\pi i\bigl(\operatorname{Res}_{z_1}(f)+\operatorname{Res}_{z_2}(f)\bigr)
=2\pi i \cdot \frac{i\,e^{-\sqrt{2}}}{4\sqrt{2}}(\sin\sqrt{2}-\cos\sqrt{2})
=\frac{-\pi e^{-\sqrt{2}}}{2\sqrt{2}}(\sin\sqrt{2}-\cos\sqrt{2}).
\]
Taking the real part (the result is already real):
\[
\boxed{\int_{-\infty}^{\infty} \frac{x^2}{16+x^4}\cos x\,dx
=\frac{\pi e^{-\sqrt{2}}}{2\sqrt{2}}\bigl(\cos\sqrt{2}-\sin\sqrt{2}\bigr)
=\frac{\pi\sqrt{2}}{4}\,e^{-\sqrt{2}}\bigl(\cos\sqrt{2}-\sin\sqrt{2}\bigr).}
\]

\section{Integrals of rational functions of trigonometric functions}

A second important class of real integrals that can be computed via residues are integrals of the form
\[
\int_0^{2\pi} \frac{P(\cos t,\sin t)}{Q(\cos t,\sin t)}\,dt,
\]
where $P,Q$ are polynomials in two variables and $Q(x,y)\neq 0$ on the unit circle $\{x^2+y^2=1\}$.

\subsection{The substitution \texorpdfstring{$z=e^{it}$}{z = exp(it)}}

The idea is to convert this real integral over $[0,2\pi]$ into a complex contour integral over the unit circle. Recall:
\[
\cos t = \frac{e^{it}+e^{-it}}{2}, \qquad \sin t = \frac{e^{it}-e^{-it}}{2i}.
\]
If we set $z=e^{it}$ then $dz=ie^{it}\,dt$, so $dt=\dfrac{dz}{iz}$, and:
\[
\cos t = \frac{1}{2}\left(z+\frac{1}{z}\right), \qquad
\sin t = \frac{1}{2i}\left(z-\frac{1}{z}\right).
\]
As $t$ runs from $0$ to $2\pi$, $z$ traces the unit circle $\gamma$ once counterclockwise. Hence
\[
\int_0^{2\pi} \frac{P(\cos t,\sin t)}{Q(\cos t,\sin t)}\,dt
=\int_\gamma f(z)\,dz,
\]
where
\[
f(z)=\frac{P\!\left(\tfrac{z+1/z}{2},\tfrac{z-1/z}{2i}\right)}{Q\!\left(\tfrac{z+1/z}{2},\tfrac{z-1/z}{2i}\right)}\cdot \frac{1}{iz}.
\]

\begin{definition}[Meromorphic function]
A function which is holomorphic on an open set $D$ except at isolated singular points (poles) is called \emph{meromorphic} on $D$.
\end{definition}

Since $P$ and $Q$ are polynomials, $f$ is meromorphic on $\mathbb{C}$. If $z_1,\dots,z_\ell$ are the poles of $f$ inside the unit disk, the residue theorem gives:
\[
\int_0^{2\pi} \frac{P(\cos t,\sin t)}{Q(\cos t,\sin t)}\,dt
=2\pi i\sum_{k=1}^\ell \operatorname{Res}_{z_k}(f).
\]

\subsection{Example}

Compute
\[
J=\int_0^{2\pi} \frac{1}{2+\cos t}\,dt.
\]
This fits the framework with $P(x,y)=1$ and $Q(x,y)=2+x$. Substituting $z=e^{it}$:
\[
2+\cos t = 2 + \frac{1}{2}\left(z+\frac{1}{z}\right) = \frac{z^2+4z+1}{2z}.
\]
Therefore
\[
f(z)=\frac{1}{\dfrac{z^2+4z+1}{2z}}\cdot \frac{1}{iz}
=\frac{2z}{z^2+4z+1}\cdot \frac{1}{iz}
=\frac{2}{i}\cdot\frac{1}{z^2+4z+1}.
\]

\paragraph{Finding the poles.}
We solve $z^2+4z+1=0$:
\[
z=\frac{-4\pm\sqrt{16-4}}{2}=-2\pm\sqrt{3}.
\]
So the poles are $z_1=-2+\sqrt{3}\approx -0.27$ and $z_2=-2-\sqrt{3}\approx -3.73$. Only $z_1$ lies inside the unit circle (since $|z_1|<1$ but $|z_2|>1$).

\paragraph{Computing the residue.}
$z_1$ is a simple pole, so:
\[
\operatorname{Res}_{z_1}(f)
=\lim_{z\to z_1}(z-z_1)f(z)
=\frac{2}{i}\cdot \frac{1}{z_1-z_2}
=\frac{2}{i}\cdot \frac{1}{2\sqrt{3}}
=\frac{1}{i\sqrt{3}}.
\]

\paragraph{Conclusion.}
\[
J=2\pi i\cdot \frac{1}{i\sqrt{3}}=\frac{2\pi}{\sqrt{3}}=\frac{2\pi\sqrt{3}}{3}.
\]

\section{General strategy}

To summarize, the global strategy for applying the residue theorem to real integrals consists in:
\begin{enumerate}
\item \textbf{Identifying the type of integral} (rational function on $\mathbb{R}$, rational function times $e^{iax}$, trigonometric integral, \dots).
\item \textbf{Extending the integrand} to a complex function $f(z)$.
\item \textbf{Choosing an appropriate closed contour} (semicircle in the upper half-plane, unit circle, \dots).
\item \textbf{Showing that the auxiliary part of the contour gives a vanishing contribution} in the appropriate limit (using estimates on $|f|$).
\item \textbf{Identifying the poles} of $f$ inside the contour and computing their residues.
\item \textbf{Applying the residue theorem} and recovering the desired real integral by taking real or imaginary parts if necessary.
\end{enumerate}

\begin{remark}
The choice of contour depends on the integrand:
\begin{itemize}
\item For $\int_{-\infty}^\infty R(x)e^{iax}\,dx$ with $a>0$: close in the upper half-plane.
\item For $a<0$: close in the lower half-plane (the orientation is then clockwise, hence a sign change).
\item For trigonometric integrals over $[0,2\pi]$: use the unit circle.
\item For other situations (integrals over $[0,\infty)$, branch cuts, etc.) one may need keyhole contours, sectors, or other shapes — but the principle remains the same.
\end{itemize}
\end{remark}
